% Source: zcash/zips commit 2bdf39494320e5e04057ac93b1794031ebfb00b5, protocol/protocol.tex lines 14384-14508
% Spec version: v2026.7.0-344-g2bdf39 [NU6.2], section 7.7.1 "Equihash"
% Pinned: 2026-10-24

\lsubsubsection{Equihash}{equihash}

An instance of the \defining{\Equihash} algorithm is parameterized by positive integers $n$ and $k$,
such that $n$ is a multiple of $k+1$. We assume $k \geq 3$.

The \Equihash parameters for \Mainnet and \Testnet are $n = 200, k = 9$.

\Equihash is based on a variation of the Generalized Birthday Problem \cite{AR2017}:
given a sequence $X_\barerange{1}{\rmN}$ of $n$-bit strings, find $2^k$ distinct
$X_{i_j}$ such that $\sxor{j=1}{2^k} X_{i_j} = 0$.

In \Equihash, $\rmN = 2^{\frac{n}{k+1}+1}$, and the sequence $X_\barerange{1}{\rmN}$ is
derived from the \blockHeader and a nonce.

\newsavebox{\powheaderbox}
\begin{lrbox}{\powheaderbox}
\begin{bytefield}[bitwidth=0.054em]{1152}
    \sbitbox{136}{$32$-bit $\nVersion$} &
    \sbitbox{256}{$256$-bit $\hashPrevBlock$} &
    \sbitbox{256}{$256$-bit $\hashMerkleRoot$} \\
    \sbitbox{256}{$256$-bit $\hashReserved$} &
    \sbitbox{136}{$32$-bit $\nTimeField$} &
    \sbitbox{136}{$32$-bit $\nBitsField$} &
    \sbitbox{256}{$256$-bit $\nNonce$}
\end{bytefield}
\end{lrbox}

\vspace{0.5ex}
Let $\powheader := \Justthebox[-5.2ex]{\powheaderbox}$

\vspace{1ex}
For $i \in \range{1}{N}$, let $X_i = \EquihashGen{n, k}(\powheader, i)$.

$\EquihashGen{}$ is instantiated in \crossref{equihashgen}.

Define $\ItoBEBSP{} \typecolon (\ell \typecolon \Nat) \times \binaryrange{\ell} \rightarrow \bitseq{\ell}$
as in \crossref{endian}.

\introsection
A \defining{\validEquihashSolution} is then a sequence $i \typecolon \range{1}{N}^{2^k}$ that
satisfies the following conditions:

\callout{}{Generalized Birthday condition}
$\vxor{j=1}{2^k} X_{i_j} = 0$.

\callout{}{Algorithm Binding conditions}
\begin{itemize}
  \item For all $r \in \range{1}{k\!-\!1}$, for all $w \in \binaryrange{k-r}:
        \smash{\vxor{j=1}{2^r}} X_{i_{w \mult \scalebox{0.65}[0.6]{$2^r$} + j}}$ has $\frac{n \mult r}{k+1}$ leading zeros; and
  \item For all $r \in \range{1}{k}$, for all $w \in \binaryrange{k-r}:
        i_{w \mult 2^r + 1 .. w \mult 2^r + 2^{r-1}} <
        i_{w \mult 2^r + 2^{r-1} + 1 .. w \mult 2^r + 2^r}$ lexicographically.
\end{itemize}

\vspace{-1ex}
\begin{pnotes}
  \item This does not include a difficulty condition, because here we are
        defining validity of an \Equihash solution independent of difficulty.
  \item Previous versions of this specification incorrectly specified the
        range of $r$ to be $\range{1}{k\!-\!1}$ for both parts of the algorithm
        binding condition. The implementation in \zcashd was as intended.
\end{pnotes}

\introlist
An \Equihash solution with $n = 200$ and $k = 9$ is encoded in the $\solution$
field of a \blockHeader as follows:

\newsavebox{\solutionbox}
\begin{lrbox}{\solutionbox}
\begin{bytefield}[bitwidth=0.45em]{105}
    \sbitbox{21}{$\ItoBEBSP{21}(i_1-1)$} &
    \sbitbox{21}{$\ItoBEBSP{21}(i_2-1)$} &
    \sbitbox{42}{$\cdots$} &
    \sbitbox{21}{$\ItoBEBSP{21}(i_{512}-1)$}
\end{bytefield}
\end{lrbox}

\newcommand{\zb}{\sbitbox{1}{$0$}}
\newcommand{\ob}{\sbitbox{1}{$1$}}
\newsavebox{\eqexamplebox}
\begin{lrbox}{\eqexamplebox}
\begin{bytefield}[bitwidth=0.75em]{63}
    \sbitbox{21}{$\ItoBEBSP{21}(68)$} &
    \sbitbox{21}{$\ItoBEBSP{21}(41)$} &
    \sbitbox{21}{$\ItoBEBSP{21}(2^{21}-1)$} \\
    \zb\zb\zb\zb\zb\zb\zb\zb\zb\zb\zb\zb\zb\zb\ob\zb\zb\zb\ob\zb\zb
    \zb\zb\zb\zb\zb\zb\zb\zb\zb\zb\zb\zb\zb\zb\zb\ob\zb\ob\zb\zb\ob
    \ob\ob\ob\ob\ob\ob\ob\ob\ob\ob\ob\ob\ob\ob\ob\ob\ob\ob\ob\ob\ob \\
    \sbitbox{8}{8-bit $0$}
    \sbitbox{8}{8-bit $2$}
    \sbitbox{8}{8-bit $32$}
    \sbitbox{8}{8-bit $0$}
    \sbitbox{8}{8-bit $10$}
    \sbitbox{8}{8-bit $127$}
    \sbitbox{8}{8-bit $255$}
    \sbitbox{7}{$\cdots$}
\end{bytefield}
\end{lrbox}

\begin{formulae}
  \item $\Justthebox{\solutionbox}$
\end{formulae}

\introlist
Recall from \crossref{bitlayout} that the bits in the above diagram are
ordered from most to least significant in each byte.
For example, if the first $3$ elements of $i$ are $[69, 42, 2^{21}]$,
then the corresponding bit array is:

\begin{formulae}
  \item $\Justthebox{\eqexamplebox}$
\end{formulae}

and so the first $7$ bytes of $\solution$ would be
$[0, 2, 32, 0, 10, 127, 255]$.

\pnote{
$\ItoBEBSP{}$ is \bigEndian, while integer field encodings in $\powheader$
and in the instantiation of $\EquihashGen{}$ are \littleEndian.
The rationale for this is that \littleEndian serialization of
\blockHeaders is consistent with \Bitcoin, but \littleEndian
ordering of bits in the solution encoding would require bit-reversal
(as opposed to only shifting).
}

